\section{Data layer}
\label{system-architecture-and-implementation:data-layer}
The benchmark runs on a single logical dataset of Norwegian building polygons, materialized as three physical representations that populate the format axis of the experimental design: GeoParquet on object storage, server-managed tables in PostGIS, and a Shapefile bundle on local disk. All three derive from the same source data under one fixed release identifier, \texttt{2026-05-23.1}, so the rows a configuration reads differ only in how they are encoded and stored, never in what they contain. Fixing the dataset across every level is the control established in Section~\ref{research-design-and-methodology:overview-and-experimental-design:factors-and-levels}: because the data are identical, a difference in transferred bytes or elapsed time traces to the format-and-engine pairing under test rather than to a difference in inputs. The subsections below cover the logical dataset and its size tiers first, then each physical representation in turn.

\subsection{Logical dataset and size tiers}
\label{system-architecture-and-implementation:data-layer:logical-dataset-and-size-tiers}
The dataset is materialized at three size tiers, small, medium, and large, that realize those defined for the workloads in Section~\ref{research-design-and-methodology:workloads:workload-definitions}. The small tier holds roughly \num{5000000} building polygons and comes directly from the test dataset pipeline, which conflates \acrfull{osm} and \acrfull{fkb} building data into a single deduplicated set. The medium and large tiers are synthesized from the small tier by cloning each polygon. Each source polygon is copied a fixed number of times. Every clone is translated to a new position within its source county, rotated by a random angle, and perturbed by a small coordinate jitter. Clones whose geometry fails a validity check are discarded. The medium tier applies nine clones per source polygon for a target of roughly \num{40000000} rows, and the large tier applies \num{23} for roughly \num{100000000}. Synthetic clones inherit their source's schema but carry \texttt{NULL} in every non-geometry attribute, so only the originals have populated values. Every row across all three tiers is assigned a \texttt{partition\_key}: a precision-3 geohash computed over the polygon centroid in the LAEA Europe projection (EPSG:3035). This key groups spatially adjacent buildings together, supporting the layout described in the following subsections.

\subsection{GeoParquet on ADLS Gen2}
\label{system-architecture-and-implementation:data-layer:geoparquet-on-adls-gen2}

The buildings dataset is written to \acrshort{adls} as GeoParquet conforming to version 1.1.0 of the specification \parencite{ogc2024_geoparquet_specification_v110}. Each partition encodes geometry as \acrshort{wkb}, adds a covering bounding-box column as a struct of \texttt{xmin}, \texttt{ymin}, \texttt{xmax}, and \texttt{ymax}, and uses a row-group size of \num{100000} rows. Columns use snappy compression, and each partition's rows are sorted by \texttt{partition\_key} before they are written. A Hive-style path keyed by release, size, theme, and region organizes the files, taking the form \texttt{\url{release/{release}/size={size}/theme=buildings/region={county}/part\_XXXXX.parquet}}. One logical release therefore resolves to a tree of per-county GeoParquet files at a fixed size tier.

This layout exposes the dataset to the block-level pruning mechanism described in Section~\ref{background:geometry-and-indexing:spatial-indexing}. The covering bounding box and the per-row-group statistics that Parquet records in the file footer let a reader compare a spatial predicate against each row group's extent and skip the groups that cannot satisfy it. This realizes the range-request access pattern established in Section~\ref{background:spatial-data-formats-and-storage-models:summary}. The region segment of the Hive path extends the same principle to whole files, so a query bounded to one county reads only the GeoParquet files written under that region rather than the national tree.

\subsection{PostGIS on Azure PostgreSQL}
\label{system-architecture-and-implementation:data-layer:postgis-on-azure-postgresql}
% One buildings_{size} table per tier (buildings_small/medium/large), GiST index
%   buildings_{size}_geometry_idx, ANALYZE after load. Loaded once by setup (streamed from blob
%   via DuckDB fetch_df_chunk to avoid materializing the full dataset). Server params from README:
%   shared_buffers=2097152 (\SI{2}{\gibi\byte}), effective_cache_size=6291456 (\SI{6}{\gibi\byte}),
%   work_mem=65536 (\SI{64}{\mebi\byte}); azure.extensions=POSTGIS. \acrfull{gist} first body use
%   if not expanded earlier. This is where the GiST/ANALYZE facts from Ch.4 threats table land.

\subsection{Shapefile bundle}
\label{system-architecture-and-implementation:data-layer:shapefile-bundle}
% Components written by benchmark_service / setup: .shp/.shx/.dbf/.cpg + .prj (WKT1_ESRI, EPSG:4326)
%   + .qix quadtree spatial index (ogr CREATE SPATIAL INDEX). Uploaded to copies/shapefile and
%   downloaded to local disk BEFORE the timed scope (laptop-workflow emulation), small tier only.
%   State the asymmetry-by-design and point back to Ch.4 threats
%   (Section~\ref{research-design-and-methodology:threats-to-validity}); do not re-argue it.